% =====================================================================
%  แง่ที่ 4 : บทความ LaTeX  (ผลลัพธ์จาก PROMPT ใหม่ / Few-shot)
%  ใช้กับ Overleaf  :  Menu -> Compiler -> XeLaTeX   (สำคัญมาก!)
%
%  วิธีเตรียมฟอนต์ไทย (ทำครั้งเดียว):
%   1. โหลด Sarabun จาก https://fonts.google.com/specimen/Sarabun (กด Get font -> Download all)
%   2. ใน Overleaf กด Upload -> New Folder ชื่อ  fonts
%   3. อัปโหลดไฟล์ 4 ไฟล์นี้เข้าโฟลเดอร์ fonts :
%        Sarabun-Regular.ttf  Sarabun-Bold.ttf
%        Sarabun-Italic.ttf   Sarabun-BoldItalic.ttf
% =====================================================================

\documentclass[12pt, a4paper]{article}

% ---------- ฟอนต์และภาษาไทย ----------
\usepackage{fontspec}
\setmainfont{Sarabun}[
  Path            = ./fonts/,
  Extension       = .ttf,
  UprightFont     = *-Regular,
  BoldFont        = *-Bold,
  ItalicFont      = *-Italic,
  BoldItalicFont  = *-BoldItalic
]
% ตัดคำภาษาไทยให้ขึ้นบรรทัดใหม่ถูกตำแหน่ง
\XeTeXlinebreaklocale "th"
\XeTeXlinebreakskip = 0pt plus 0.1pt

% ---------- แพ็กเกจ ----------
\usepackage[margin=2.5cm]{geometry}
\usepackage{amsmath, amssymb}
\usepackage{booktabs}
\usepackage{graphicx}
\usepackage{float}
\usepackage{caption}
\usepackage{enumitem}
\usepackage[hidelinks]{hyperref}
\usepackage{xcolor}

\definecolor{accent}{HTML}{065A82}
\captionsetup{font=small, labelfont={bf,color=accent}}
\setlength{\parindent}{1.5em}
\setlength{\parskip}{0.4em}

% ---------- ข้อมูลหัวเรื่อง ----------
\title{\textbf{การเรียนรู้ของเครื่อง:\\ จากหลักการทางคณิตศาสตร์สู่การนำไปใช้งานจริง}}
\author{ชื่อ--นามสกุล \\ รหัสนักศึกษา \\ รายวิชา ปัญญาประดิษฐ์}
\date{\today}

\begin{document}
\maketitle

% =====================================================================
\begin{abstract}
\noindent
บทความนี้นำเสนอภาพรวมของการเรียนรู้ของเครื่อง (Machine Learning) โดยเริ่มจากการนิยามปัญหาในเชิงคณิตศาสตร์
อธิบายบทบาทของฟังก์ชันกระตุ้น (activation function) และฟังก์ชันสูญเสีย (loss function)
ตามด้วยขั้นตอนวิธีการลดตามความชัน (gradient descent) ซึ่งเป็นกลไกหลักในการปรับค่าพารามิเตอร์
ส่วนท้ายเปรียบเทียบอัลกอริทึมพื้นฐานสามชนิดและอภิปรายปัญหาการเรียนรู้เกินพอดี (overfitting)
พร้อมแนวทางบรรเทา บทความนี้จัดทำขึ้นเป็นส่วนหนึ่งของการฝึกพัฒนา prompt
เพื่อสร้างเอกสารวิชาการด้วยระบบเรียงพิมพ์ \LaTeX{}

\vspace{0.5em}
\noindent\textbf{คำสำคัญ:} การเรียนรู้ของเครื่อง, การลดตามความชัน, ฟังก์ชันสูญเสีย, การเรียนรู้เกินพอดี
\end{abstract}

% =====================================================================
\section{บทนำ}

การเรียนรู้ของเครื่องคือสาขาหนึ่งของปัญญาประดิษฐ์ที่มุ่งสร้างระบบซึ่งปรับปรุงประสิทธิภาพของตนเองได้
จากประสบการณ์ แทนที่จะเขียนกฎเกณฑ์ทุกข้อลงไปตรง ๆ ผู้พัฒนาเพียงกำหนดรูปแบบของฟังก์ชัน
แล้วให้ข้อมูลเป็นตัวกำหนดค่าพารามิเตอร์ที่เหมาะสม

ในเชิงรูปนัย เมื่อกำหนดชุดข้อมูลฝึก
\[
  \mathcal{D} = \{(\mathbf{x}^{(i)}, y^{(i)})\}_{i=1}^{n},
  \qquad \mathbf{x}^{(i)} \in \mathbb{R}^{d},\ y^{(i)} \in \mathcal{Y}
\]
เป้าหมายคือการค้นหาฟังก์ชัน $f_{\theta} : \mathbb{R}^{d} \to \mathcal{Y}$
ที่ทำให้ค่าความสูญเสียเฉลี่ยต่ำที่สุด

\begin{equation}
  \theta^{*} \;=\; \arg\min_{\theta}\;
  \frac{1}{n}\sum_{i=1}^{n} \mathcal{L}\!\left(f_{\theta}(\mathbf{x}^{(i)}),\, y^{(i)}\right)
  \;+\; \lambda\,\Omega(\theta)
  \label{eq:erm}
\end{equation}

โดยพจน์ $\Omega(\theta)$ คือพจน์ควบคุมความซับซ้อน และ $\lambda \ge 0$ คือค่าถ่วงน้ำหนัก

% =====================================================================
\section{ฟังก์ชันกระตุ้น}

ฟังก์ชันกระตุ้นทำหน้าที่เพิ่มความไม่เป็นเชิงเส้นให้กับโครงข่ายประสาทเทียม
หากปราศจากฟังก์ชันเหล่านี้ โครงข่ายหลายชั้นจะยุบรวมเป็นการแปลงเชิงเส้นชั้นเดียวเสมอ
ฟังก์ชันที่ใช้บ่อยมีสามรูปแบบดังนี้

\begin{align}
  \sigma(z) &= \frac{1}{1 + e^{-z}} & &\in (0,1) \label{eq:sigmoid}\\[4pt]
  \tanh(z)  &= \frac{e^{z} - e^{-z}}{e^{z} + e^{-z}} & &\in (-1,1) \label{eq:tanh}\\[4pt]
  \mathrm{ReLU}(z) &= \max(0, z) & &\in [0,\infty) \label{eq:relu}
\end{align}

อนุพันธ์ของฟังก์ชันซิกมอยด์ในสมการ~\eqref{eq:sigmoid} มีรูปแบบที่กระชับเป็นพิเศษ

\begin{equation}
  \frac{d\sigma}{dz} = \sigma(z)\bigl(1 - \sigma(z)\bigr)
  \label{eq:sigmoid-deriv}
\end{equation}

จากสมการ~\eqref{eq:sigmoid-deriv} จะเห็นว่าเมื่อ $|z|$ มีค่ามาก อนุพันธ์จะเข้าใกล้ศูนย์
ทำให้สัญญาณความชันที่ส่งย้อนกลับลดลงจนแทบหยุดการเรียนรู้ ปรากฏการณ์นี้เรียกว่า
\emph{vanishing gradient} และเป็นเหตุผลสำคัญที่ทำให้ ReLU ได้รับความนิยมในโครงข่ายลึก

% =====================================================================
\section{การลดตามความชัน}

การหาคำตอบของสมการ~\eqref{eq:erm} โดยตรงมักทำไม่ได้ในทางปฏิบัติ
จึงใช้วิธีการทำซ้ำเพื่อเคลื่อนค่าพารามิเตอร์ไปในทิศทางที่ค่าความสูญเสียลดลง

\begin{equation}
  \theta_{t+1} \;=\; \theta_{t} \;-\; \eta \, \nabla_{\theta} \mathcal{L}(\theta_{t})
  \label{eq:gd}
\end{equation}

ค่า $\eta$ เรียกว่าอัตราการเรียนรู้ (learning rate) และเป็นไฮเปอร์พารามิเตอร์ที่อ่อนไหวที่สุดตัวหนึ่ง

\begin{itemize}[leftmargin=2em, itemsep=0.2em]
  \item ถ้า $\eta$ เล็กเกินไป การลู่เข้าจะช้ามากและอาจติดอยู่ที่ค่าต่ำสุดเฉพาะที่
  \item ถ้า $\eta$ ใหญ่เกินไป ค่าพารามิเตอร์จะแกว่งข้ามจุดต่ำสุดและอาจลู่ออก
\end{itemize}

% =====================================================================
\section{การเปรียบเทียบอัลกอริทึมพื้นฐาน}

ตารางที่~\ref{tab:compare} เปรียบเทียบอัลกอริทึมสามชนิดที่ใช้บ่อยในงานจำแนกประเภท

\begin{table}[H]
\centering
\caption{การเปรียบเทียบอัลกอริทึมการเรียนรู้ของเครื่องเชิงกำกับ}
\label{tab:compare}
\small
\begin{tabular}{@{}llll@{}}
\toprule
\textbf{คุณลักษณะ} & \textbf{Logistic Regression} & \textbf{Decision Tree} & \textbf{Neural Network} \\
\midrule
รูปแบบขอบเขตตัดสิน & เชิงเส้น         & ขั้นบันได         & ไม่เป็นเชิงเส้นอิสระ \\
ความสามารถอธิบายผล & สูง               & สูงมาก            & ต่ำ \\
ปริมาณข้อมูลที่ต้องการ & น้อย           & ปานกลาง           & มาก \\
เวลาที่ใช้ฝึก        & เร็ว              & เร็ว               & ช้า \\
ความเสี่ยง overfitting & ต่ำ            & สูง                & สูง \\
เหมาะกับงาน          & ข้อมูลตาราง       & ข้อมูลตาราง        & ภาพ เสียง ข้อความ \\
\bottomrule
\end{tabular}
\end{table}

% =====================================================================
\section{การเรียนรู้เกินพอดี}

การเรียนรู้เกินพอดีเกิดขึ้นเมื่อโมเดลจดจำรายละเอียดและสัญญาณรบกวนของชุดข้อมูลฝึก
จนสูญเสียความสามารถในการวางนัยทั่วไป สังเกตได้จากค่าความสูญเสียบนชุดฝึกที่ลดลงเรื่อย ๆ
สวนทางกับค่าความสูญเสียบนชุดตรวจสอบที่กลับเพิ่มขึ้น

แนวทางบรรเทาที่ใช้ได้ผลมีดังนี้

\begin{enumerate}[leftmargin=2em, itemsep=0.2em]
  \item เพิ่มปริมาณข้อมูลฝึก หรือใช้เทคนิคเพิ่มข้อมูลสังเคราะห์ (data augmentation)
  \item เพิ่มพจน์ควบคุมความซับซ้อน เช่น $L_{2}$ ซึ่งกำหนดให้ $\Omega(\theta) = \lVert \theta \rVert_{2}^{2}$
  \item หยุดการฝึกก่อนกำหนด (early stopping) เมื่อค่าความสูญเสียบนชุดตรวจสอบเริ่มเพิ่มขึ้น
  \item ใช้ dropout เพื่อสุ่มปิดหน่วยประมวลผลบางส่วนระหว่างการฝึก
\end{enumerate}

% =====================================================================
\section{บทสรุป}

บทความนี้ได้ทบทวนองค์ประกอบหลักของการเรียนรู้ของเครื่อง ตั้งแต่การนิยามปัญหาการหาค่าเหมาะที่สุด
ในสมการ~\eqref{eq:erm} บทบาทของฟังก์ชันกระตุ้น กลไกการปรับพารามิเตอร์ด้วยสมการ~\eqref{eq:gd}
ไปจนถึงการเปรียบเทียบอัลกอริทึมและการรับมือกับการเรียนรู้เกินพอดี
ประเด็นสำคัญที่ได้จากการทบทวนคือ คุณภาพของแบบจำลองขึ้นอยู่กับคุณภาพของข้อมูล
และความเหมาะสมของสมมติฐานที่ใส่เข้าไป มากกว่าความซับซ้อนของอัลกอริทึมเพียงอย่างเดียว

% =====================================================================
\begin{thebibliography}{9}
\bibitem{goodfellow}
  Goodfellow, I., Bengio, Y., and Courville, A.
  \emph{Deep Learning}. MIT Press, 2016.

\bibitem{bishop}
  Bishop, C. M.
  \emph{Pattern Recognition and Machine Learning}. Springer, 2006.

\bibitem{hastie}
  Hastie, T., Tibshirani, R., and Friedman, J.
  \emph{The Elements of Statistical Learning}. 2nd ed., Springer, 2009.
\end{thebibliography}

\end{document}
